\chapter{Introduction}\label{ch:intro}

\section{Motivation}

\begin{itemize}
\item Epilepsy affected an estimated 45.9 million people in 2016 \cite{gbd2019epilepsy}. Patients often miss their own
  seizures, so diaries undercount them \cite{elger2018underreporting}, and tonic-clonic seizures carry a risk of sudden
  death \cite{ryvlin2013mortemus}.
\item Wearables are recommended today only for tonic-clonic seizures \cite{beniczky2021guideline}; for focal seizures
  (822 of the 883 in SeizeIT2) false alarms are the main obstacle \cite{swinnen2026seizeit2}.
\item Behind-the-ear EEG with ECG, EMG and motion is a discreet alternative \cite{vandecasteele2020bte,vandecasteele2021ecg};
  SeizeIT2 \cite{bhagubai2025seizeit2} provides 125 patients recorded with such a wearable.
\item A wearable detector must run for days on a small battery. Multiplications dominate the cost of neural networks
  in hardware, while an FPGA's basic element is a 6-input lookup table (LUT). \emph{LUT-native} networks
  \cite{bacellar2024dwn,petersen2022dlgn,umuroglu2020logicnets} are built from exactly this element.
\item A medical device must also not fail silently: a hardware fault in the classifier turns into missed seizures or
  false alarms, and an alarm is only useful if the right person can act on it quickly.
\end{itemize}

\section{Problem Statement and Objectives}

\begin{quote}
\itshape Can a LUT-native network, fed by integer features from behind-the-ear EEG, ECG, EMG and motion, detect
seizures in unseen patients as well as the published SeizeIT2 baselines, as multiplier-free hardware that can test
itself, run on a real FPGA board, and alert a doctor?
\end{quote}

\begin{table}[H]
\centering\small
\caption{Objectives and their status.}\label{tab:objectives}
\begin{tabularx}{\textwidth}{@{}cLP{42mm}@{}}
\toprule
\textbf{\#} & \textbf{Objective} & \textbf{Status}\\
\midrule
1 & Audit SeizeIT2 and reproduce its published baselines & Done\\
2 & Integer-only, bit-exact multimodal feature pipeline & Done\\
3 & Compare LUT-native, tree and multiplier-based models on accuracy and hardware cost & Done\\
4 & Extend the LUT network for wearables (distillation, missing sensors, post-ictal verification) & Done\\
5 & Evaluate the frozen design once on the 29 held-out test patients & Done\\
6 & Self-testing: complete single-fault analysis and an on-chip test that detects every relevant fault & Done, measured on the board\\
7 & Front end and classifier in the FPGA fabric; KV260 implementation & Done, all 9 board tests pass\\
8 & Live demo with alarm LED and buzzer & Done on the board\\
9 & Doctor alert through an ESP32 (phone notification and live page) & Built and tested in emulation; ESP32 hardware test pending\\
\bottomrule
\end{tabularx}
\end{table}

\section{Contributions}

\begin{itemize}
\item The first LUT-native seizure detector, and the first benchmark of LUT-family models (DWN, logic-gate network,
  TreeLUT) against standard machine learning on SeizeIT2 with event-level clinical scoring.
\item A bit-exact, multiplier-free integer front end, and a microcoded engine that computes it in the FPGA fabric:
  the whole detector is 1,202 LUTs with no multiplier.
\item Distillation from a tree ensemble into the LUT network (false alarms per day from 2.4 to 0.9) and a LUT-native
  post-ictal verification stage that detects more seizures than the dataset paper's SVM.
\item \emph{Self-testing}: every one of the 3,592 single faults of the classifier judged by its effect on decisions,
  and 255 selected test vectors that detect all decision-relevant faults; verified fault by fault on the board.
\item A complete Kria KV260 implementation with a nine-step board test campaign, a live demo, and a doctor alert
  that tells the doctor what to do.
\item An independent reproduction of the SeizeIT2 ChronoNet baseline, including a protocol with no test information.
\end{itemize}

\section{Team}

\begin{table}[H]
\centering\small
\caption{Project team.}\label{tab:team}
\begin{tabular}{@{}ll@{}}
\toprule
\textbf{Name} & \textbf{Roll number}\\
\midrule
\studentA & \rollA\\
\studentB & \rollB\\
\studentC & \rollC\\
\bottomrule
\end{tabular}
\end{table}

\section{Organisation of the Report}

Chapter~\ref{ch:lit} reviews the literature, Chapter~\ref{ch:method} the dataset and method, Chapter~\ref{ch:results}
the detection results, Chapter~\ref{ch:hw} the FPGA hardware, self-test and board results, Chapter~\ref{ch:alert} the
doctor alert, Chapter~\ref{ch:challenges} the challenges, and Chapter~\ref{ch:concl} the conclusions, the publication
plan and future work.
